\subsubsection{Results and Impact}\label{sec:results-and-impact}

We have talked about the limitations of N-gram models, Bengio's Neural Language Model, and the architecture and training process. The natural question now is: \emph{\textbf{did it actually work?}} The answer is \textbf{yes}. Bengio's model demonstrated that neural networks could outperform traditional N-gram language models while simultaneously learning meaningful word representations. This was a turning point in Natural Language Processing.

\highspace
\begin{flushleft}
    \textcolor{Green3}{\faIcon{tools} \textbf{Experimental Results}}
\end{flushleft}
The Bengio et al. (2003) paper evaluated their Neural Language Model on several benchmark datasets, including the Brown Corpus and the AP News Corpus.

\highspace
\textcolor{Green3}{\faIcon{tools} \textbf{Network Configurations.}} The paper used relatively small embedding dimensions compared to modern standards.
\begin{itemize}
    \item \textbf{Brown Corpus}
    \begin{equation*}
        h = 100, \quad n = 5, \quad m = 30
    \end{equation*}
    \item \textbf{AP News Corpus}
    \begin{equation*}
        h = 60, \quad n = 6, \quad m = 100
    \end{equation*}
\end{itemize}
Where $h$ is the hidden layer size, $n$ is the N-gram context size, and $m$ is the embedding dimension.

\highspace
\textcolor{Red2}{\faIcon{exclamation-triangle} \textbf{Training Cost.}} In 2003, training a neural language model was computationally expensive. The paper reported that on AP News Corpus, training took 3 weeks on 40 CPUs (running only 5 epochs!).\cite{bengio2003neural} Today a comparable experiment can often be run in hours or minutes. However, this \hl{highlights how computationally expensive neural NLP was in the early 2000s}, which was a significant barrier to adoption.

\highspace
\textcolor{Green3}{\faIcon{check-circle} \textbf{Performance Improvements.}} Bengio proposed a model that worked significantly better, but was computationally expensive. The paper reported that their Neural Language Model achieved \hl{24\% lower perplexity} on the Brown Corpus and \hl{8\% lower perplexity} on the AP News Corpus compared to traditional N-gram models.

\begin{deepeningbox}[: What is perplexity?]
    Perplexity is a common evaluation metric for language models. Imagine a multiple-choice exam where we have to choose a word from a set of options to complete a sentence. For example, given the sentence ``\emph{The cat is \_\_\_.}'', and the options are ``\emph{sleeping}'', ``\emph{driving}'', ``\emph{flying}'', and ``\emph{quantum}''. We immediately think ``\emph{sleeping}'' is the most likely word to complete the sentence, and our uncertainty is low. Instead, if the options were ``\emph{sleeping}'', ``\emph{eating}'', ``\emph{running}'', and ``\emph{hunting}'' we might be less certain about which word is correct, and our uncertainty is higher. \hl{Perplexity measures this uncertainty}.

    \highspace
    A language model is constantly asking ``\emph{what is the next word?}'', and \definition{Perplexity} measures ``\emph{how confused is the model?}'' when trying to predict the next word, or similarly, ``\emph{how many plausible next words am I considering?}''. \hl{Higher perplexity} means the model has a \hl{larger set of plausible next words}, while \hl{lower perplexity} means the \hl{model is more confident in its prediction.} Therefore, a lower perplexity indicates a better language model.

    \highspace
    \textcolor{Green3}{\faIcon{question-circle} \textbf{Why Bengio's result was impressive?}} Suppose an N-gram model had a perplexity of 200 on a dataset. A 24\% reduction in perplexity means the Neural Language Model would have a perplexity of 152, thus \hl{the model understands language structure better and can predict the next word more confidently.}
\end{deepeningbox}

\highspace
\textcolor{Green3}{\faIcon{book} \textbf{Generalization Capability.}} The most important contribution of Bengio's Neural Language Model was not just the performance improvement, but the fact that it could generalize to unseen N-grams. Traditional N-gram models would assign zero probability to any N-gram not seen during training, while the Neural Language Model could still make reasonable predictions based on learned word embeddings and shared parameters. In one word, \textbf{generalization} was the key breakthrough.

\highspace
Furthermore, one of the most fascinating aspects of Bengio's work is that \textbf{nobody explicitly told the network to learn semantic relationships between words}. The objective was simply to \emph{predict the next word}. However, the network learned to represent words in a way that captured their meanings and relationships. For example, the model could learn that ``\emph{king}'' and ``\emph{queen}'' are related, and that ``\emph{king}'' is to ``\emph{man} as ``\emph{queen}'' is to ``\emph{woman}''. This \hl{emergent property of learning meaningful word representations was a significant milestone in NLP.}

\highspace
\begin{flushleft}
    \textcolor{Red2}{\faIcon{exclamation-triangle} \textbf{Limitations of the Neural Network Language Model}}
\end{flushleft}
Despite its success, Bengio's model had significant weaknesses.
\begin{itemize}
    \item \textbf{Computation Cost}. Training was extremely slow. The giant Softmax over the vocabulary was expensive.
    \item \textbf{Scalability}. The model struggled on very large datasets and vocabularies, which limited its applicability in real-world scenarios.
    \item \textbf{Rare Words}. Words appearing only a few times receive very little training.
\end{itemize}
However, it introduced the idea of learning distributed word representations. It focused on language-model accuracy, but the learned embeddings became the most influential contribution. 